\section{Low-Level Motion Generator}

\begin{figure*}[t]
    \centering
\includegraphics[width=\textwidth]{figures/3_coarse2fine_v7_compressed.pdf}
\vspace{-6mm}
    \caption{The interaction module consists of four stages. Starting with the initial human and object poses, waypoints, final object pose, and a text instruction, the first stage generates coarse object and human motion. The relative wrist pose from this motion is then used to generate a valid grasp pose. Using this as a condition, the next stage re-generates the motion to align with the grasp constraint. Finally, the finger motion generation module produces finger motions during the approach and release of objects.} 
    \vspace{-4mm}
    \label{fig:coarse_to_fine}
\end{figure*}

Given the waypoints and textual instructions generated by the high-level planner, the low-level motion generator executes an interaction module and a navigation module sequentially to generate long-horizon human-object interaction motions.

\subsection{Background: Conditional Diffusion Model}
\label{sec:cond_diff}
The diffusion model involves a forward diffusion process and a reverse diffusion process. Starting from clean data $\bm{x}_0$,  noise is progressively introduced until, after N steps, the data becomes approximately $\bm{x}_{N} \sim \mathcal{N}(0, \bm{I})$. The denoising network is trained to reverse this process so that sampling a noise $\bm{x}_{N}$ and applying N denoising steps yields the clean data representation $\hat{\bm{x}}_0$. Each denoising step is defined as 
\begin{equation}
p_{\theta}(\bm{x}_{n-1}|\bm{x}_{n}, \bm{c}) := \mathcal{N}(\bm{x}_{n-1}; \bm{\mu}_{\theta}(\bm{x}_n, n, \bm{c}), \sigma_{n}^{2}\bm{I}),  
\end{equation}
where $p_{\theta}$ represents a neural network, $\bm{c}$ represents the conditions, $\bm{\mu}_{\theta}(\bm{x}_n, n, \bm{c})$ represents the learned mean, $\sigma_{n}^{2}$ denotes a fixed variance. The training loss is defined as $\mathcal{L} =  \mathbb{E}_{\bm{x}_0, n}||\hat{\bm{x}}_{\theta}(\bm{x}_{n}, n, \bm{c}) - \bm{x}_{0}||_{1}$.
% The detailed architecture is provided in the supplementary material. 

\subsection{Interaction Module}
\label{sec:low_level_interaction}
We introduce a multi-stage interaction module that synthesizes synchronized full-body, finger, and object motion, as shown in~\cref{fig:coarse_to_fine}.
The first stage, \textit{CoarseNet}, generates initial human and object motions without detailed finger movement.
Recognizing that finger poses typically remain static during interaction, the second stage generates a \textit{grasp} pose based on the initial results.
Then \textit{RefineNet} re-generates the human and object motions to seamlessly align with the optimized grasp pose.
Finally, \textit{FingerNet} generates smooth finger motions for approaching and releasing the object, completing the interaction sequence.

\subsubsection{CoarseNet}
\label{sec:coarse_net}
Given the initial human and object poses, waypoints, and final object pose, along with a text instruction, CoarseNet generates synchronized human motion and object motion. This is achieved using a pre-trained diffusion model CHOIS~\cite{li2023controllable}.
% We utilize the pre-trained CHOIS model to generate synchronized object and human motions based on the initial state, end object state, sparse waypoints, object geometry, and text.

% \paragraph{Motion Representation}
% We denote human motion as $\bm{H} \in \mathbb{R}^{T \times D}$, where $T$ represents the sequence length and $D$ represents the dimension of the human pose. Each $\bm{H}_{t}$ includes global joint positions and 6D continuous rotations~\cite{zhou2019continuity}, excluding finger joints. Object motion $\bm{O} \in \mathbb{R}^{T \times 12}$ includes the object's global position and relative rotation matrix with respect to the input object geometry. The model also predicts contact labels $\bm{L} \in \mathbb{R}^2$ for both hands, which are used for test-time guidance. 

% \paragraph{Object Geometry Representation. } Object geometry $\bm{G} \in \mathbb{R}^{1024 \times 3}$ is represented using the Basis Point Set (BPS)~\cite{prokudin2019efficient} method. This method involves sampling $1024$ basis points from the volume of a ball with a $1$-meter radius centered on the object. Directional vectors are then calculated from each basis point to the nearest point on the object, forming a compact BPS representation. 

% \paragraph{Input Condition Representation. } 
% The model uses a masked motion data representation $\bm{S} \in \mathbb{R}^{T \times (D+12)}$ to encode the initial states, end object state, and sparse waypoints. The initial state includes the human pose and object pose in the first frame, while the sparse waypoints provide a series of 2D object positions every 30 frames. Any remaining spaces in $\bm{S}$ are padded with zeros. The masked motion data $\bm{S}$, the BPS representation $\bm{G}$, and text embeddings extracted by a pretrained CLIP model~\cite{radford2021learning} together form the condition $\bm{c}$ for the conditional diffusion model.

% The conditional diffusion model uses input conditions $\bm{c}$ to generate object motion $\bm{O}$, human motion $\bm{H}$, and hand contact $\bm{L}$. To ensure realistic hand-object contact and reduce artifacts, CHOIS employs test-time guidance~\cite{ho2022video} to minimize wrist-object distance in the contact phase, feet-floor distance, and penalize object-floor penetration. For more details, please refer to~\cite{li2023controllable}.


\paragraph{Data Representation.} 
We denote human motion as $\bm{H} \in \mathbb{R}^{T \times D}$, where $T$ represents the sequence length and $D$ represents the dimension of the human pose. Each $\bm{H}_{t}$ includes global joint positions and 6D rotations~\cite{zhou2019continuity}, excluding finger joints. Object motion $\bm{O} \in \mathbb{R}^{T \times 12}$ includes the object's global position and rotation matrix.
CoarseNet predicts human motion, object motion, and contact labels $\bm{L} \in \mathbb{R}^2$ for both hands, resulting in $\bm{x} = \left\{\bm{H}, \bm{O}, \bm{L}\right\}$.

The input condition is represented as $\bm{c} = \left\{\bm{S}, \bm{G}, \bm{T}\right\}$. $\bm{S} \in \mathbb{R}^{T \times (D+12)}$ is a masked motion data representation. It includes the human and object pose in the first frame, the object pose in the last frame, and 2D waypoints sampled every 30 frames. Any remaining entries in $\bm{S}$ are padded with zeros. The object geometry $\bm{G}$ represents the BPS representation~\cite{prokudin2019efficient} to encode object geometry, and $\bm{T}$ refers to the text embeddings extracted using CLIP~\cite{radford2021learning}.

\vspace{-3mm}
\paragraph{Segmenting Contact Phases.} 
% Since the object is stationary at the beginning and end, 
We segment the generated interaction sequence into \textit{pre-contact}, \textit{contact}, and \textit{post-contact} phases based on the predicted contact labels $\bm{L}$.  
% The contact phase is computed separately for each hand. 
% When both hands contact the object, the contact phase is the union of left-hand and right-hand contact phases.
The contact phase is determined separately for each hand, dividing the model output $\bm{x}_\text{coarse}$ into $\bm{x}_\text{pre-contact} \oplus \bm{x}_\text{contact} \oplus \bm{x}_\text{post-contact}$.
% From $\bm{x}_\text{contact}$, we calculate the average wrist pose in the object frame, $\bm{w}$, which is then used as the input for the second stage.
We then compute the average wrist pose $\bm{w}$ from $\bm{x}_\text{contact}$, serving as input for the next stage. 

\subsubsection{Grasp Pose Generation}
\label{subsubsec:grasp_pose}
% The output of CoarseNet provides coordinated full-body and object motion, but the hands maintain a rest pose throughout the motion, resulting in visually and physically implausible hand-object contacts. We address this issue in the second stage by applying a state-of-the-art grasp pose generation method~\cite{wang2023dexgraspnet} to obtain a detailed grasp pose with accurate finger-object contact.

% Based on the object's velocity in the first stage, we segment the generated interaction sequence into pre-contact, contact, and post-contact phases, where the object velocities in the pre-contact and post-contact phases are approximately zero. 
% We also use the average relative wrist pose computed from Stage 1 to generate a plausible grasp pose as an initial guess to be optimized.
% Starting with an object mesh and the initial hand pose $\bm{g} = (\bm{T}, \bm{R}, \bm{\theta})$, where $\bm{T} \in \mathbb{R}^3$ and $\bm{R} \in SO(3)$ represent the wrist position and rotation in the object's frame, and $\bm{\theta} \in \mathbb{R}^{45}$ represents the finger configuration, we formulate an optimization to compute a hand pose that minimize a differentiable force closure energy term and a regularization term (details in \cite{wang2023dexgraspnet}). For motions involving both hands, such as carrying a box, we omit the force closure term because the grasp stability does not rely solely on the fingers of one hand. The optimized grasp pose is denoted by $\hat{\bm{g}} = (\bm{\hat{T}}, \bm{\hat{R}}, \bm{\hat{\theta}})$. The optimized wrist pose $\hat{\bm{w}} = (\bm{\hat{T}}, \bm{\hat{R}})$ is subsequently used as the input condition for the third stage.

The second stage generates the finger pose for the contact phase using a state-of-the-art grasp pose generation method~\cite{wang2023dexgraspnet}.
% Although originally designed for robotic hand grasping, it relies solely on kinematic chain information, making it compatible with SMPL-X~\cite{smplx}.
Given an object, wrist pose $\bm{w}$ and rest finger pose, an optimization is conducted to generate a physically plausible grasp pose.  
We modified \cite{wang2023dexgraspnet} to fit our needs by: (1) increasing the object surface sample points from 2,000 to 20,000 to maintain sample density for larger objects and (2) omitting the force closure term for tasks involving both hands, such as carrying a box, where stability doesn’t solely rely on a single hand.
% We made the following modifications to adapt~\cite{wang2023dexgraspnet} for our task:
% We made several modifications to tailor the optimization method to our specific application:
% \begin{itemize} 
% \item DexGraspNet~\cite{wang2023dexgraspnet} randomly samples 2,000 points on the object surface to measure hand-object penetration. Since our objects are much larger, we increase this to 20,000 sample points to maintain a similar sample density. 
% \item DexGraspNet~\cite{wang2023dexgraspnet} uses a force closure term in the optimization process to measure grasp stability. However, our tasks involve motions with both hands, such as carrying a box, where stability is not solely dependent on one hand. We omit this term for such motions.
% \end{itemize}

The optimized grasp pose is denoted by $\hat{\bm{g}} = (\hat{\bm{w}}, {\hat{\bm{\theta}}})$, where $\hat{\bm{w}}$ is the optimized wrist pose, and $\hat{\bm{\theta}}$ is the finger pose. This grasp pose $\hat{\bm{g}}$ is maintained throughout the contact phase. However, as CoarseNet’s generated motion may not fully align with $\hat{\bm{w}}$, we introduce RefineNet to address this misalignment.
% \jiaman{This is achieved by RefineNet, right?}\zhen{Yes, the wrist pose condition is achieved by RefineNet. But here I also want to stress that the finger pose should stay to $\hat{\bm{\theta}}$ during the contact phase.} \jiaman{I see. So basically here, you want to say we use the optimized grasp pose for all the timesteps in the contact phase? Then we should explicitly say this.}\zhen{Fixed}

\subsubsection{RefineNet}
\label{sec:refine_net}

% Specifically, we encode the wrist pose and contact phase information by modifying the input masked data conditions (introduced in Section 5.1). We expand the data dimension to include an additional term representing the relative wrist pose. For the frames in the contact phase, the optimized wrist pose $\hat{\bm{w}} \in \mathbb{R}^{T\times18}$ includes the relative position and 6D orientation of both wrists. For the frames in the pre-contact and post-contact phases, we pad this term with zeros.

The goal of RefineNet is to re-generate human-object interaction motions that align with the optimized grasp pose $\hat{\bm{g}}$. We adopt a conditional diffusion model based on CoarseNet, adding two conditions: \textit{wrist-object relative pose} to align the wrist with $\hat{\bm{w}}$ and \textit{object static} to keep the object stationary during non-contact phases. RefineNet’s input condition, $\bm{c}_{r} = \left\{\bm{W}, \bm{S}_{r}, \bm{G}, \bm{T}\right\}$, includes $\bm{G}$ and $\bm{T}$ from CoarseNet, with $\bm{W}$ and $\bm{S}_{r}$ detailed as follows. 


\vspace{-3mm}
\paragraph{Wrist-Object Relative Pose Condition.} 
The condition $\bm{W} \in \mathbb{R}^{T \times 18}$ represents the position and 6D rotation of both wrists in the object's frame, including only contact-phase wrist poses, with other entries set to zero.
% During training, it only includes the wrist pose for the contact phase, while all other entries are masked with zeros.
% During sampling, we set the entries corresponding to the contact phase with $\hat{\bm{w}}$, and all other entries remain zero.
% \jiaman{we can say contact phase?}\zhen{Here I’m referring to the training process. “contact phase” is the term specific to sampling, so I tried to avoid using it in this context. What do you think?} \jiaman{But there's no difference when you compute contact phases in training and sampling? Maybe we can just use contact phase. It's a bit confusing now.}\zhen{Fixed}

\vspace{-2mm}
\paragraph{Object Static Condition.}
In CoarseNet, the masked motion condition $\bm{S}$ was zero-padded except at the start, end, and waypoints. RefineNet uses contact phase information from CoarseNet to adjust $\bm{S}$ to $\bm{S_{r}}$, assigning a static object pose in pre-contact and post-contact frames.
% We apply the same approach during sampling.
% \jiaman{In training, you have GT data. GT data should already have static object poses in pre-contact and post-contact phases? Why do we need to adjust the object pose?}\zhen{Because in CoarseNet (\cref{sec:coarse_net}) training, $\bm{S}$ only contains the initial and end object pose. All other entries, including pre-contact and post-contact phases, are masked with zero. In RefineNet training, we remove the mask on the pre-contact and post-contact parts, so that the model knows how to deal with this static condition during sampling.}.\jiaman{Got it. But it seems training and sampling are the same? so we don't need to repeat this sentence. }\zhen{Fixed}

% Now, with $\bm{c} = \left\{\bm{S}, \bm{G}, \bm{T}, \bm{W}\right\}$, we further introduce a loss term to enforce the wrist-object relative pose condition.

\vspace{-2mm}
\paragraph{Wrist-Object Relative Pose Loss.}
We incorporate an additional loss to enforce the wrist-object relative pose condition. We sample 100 points $\bm{K}_{\text{rest}} \in \mathbb{R}^{100 \times 3}$ uniformly on the rest object’s surface and pre-calculate their positions in both wrists' coordinate frames, denoted as $\bm{K}_{\text{w}} \in \mathbb{R}^{2\times T \times 100 \times 3}$. 
At each time step, we compute the points' positions in the wrist frame, $\hat{\bm{K}}_{\text{w}}$, using the predicted orientations and positions of the object and wrist, $\bm{R}_{\text{o}}, \bm{T}_{\text{o}}, \bm{R}_{\text{w}}, \bm{T}_{\text{w}}$:
\begin{align}
\bm{K}_{\text{global}} &= \bm{R}_{\text{o}}\bm{K}_{\text{rest}} + \bm{T}_{\text{o}}, \\
\hat{\bm{K}}_{\text{w}} &= \bm{R}_{\text{w}}^{-1} \left( \bm{K}_{\text{global}} - \bm{T}_{\text{w}}\right).
\end{align}
The loss is computed as:
\begin{align}
\mathcal{L}_{\text{relative }}=
\sum_{t=1}^T
\bm{L}_{t}
\left\| 
\hat{\bm{K}}_{\text{w}, t}-\bm{K}_{\text{w}, t}
\right\|_1,
\end{align}
where $\bm{L}_{t}$ represents the contact labels, masking out the loss when the hand and object are not in contact.

% In addition, we adopt foot contact loss~\cite{tevet2022human} to mitigate foot sliding.



\vspace{-2mm}
\paragraph{Post-Processing.}
RefineNet generates human-object motions that align closely with input conditions, though minor adjustments may be required. To refine alignment, we apply post-processing: replacing object poses in pre- and post-contact phases with static poses and recalculating wrist trajectories during the contact phase based on object motion and the optimized wrist pose $\hat{\bm{w}}$. Additional implementation details are in the supplemental materials.






